% !TEX root = ../jalase01.tex
%% ============================================================
%%  جلسهٔ ۱ — مفاهیم اولیهٔ احتمال، نظریهٔ مجموعه‌ها، ترکیبیات
%% ============================================================
\chapter{مفاهیم اولیهٔ احتمال، نظریهٔ مجموعه‌ها و ترکیبیات}
\label{ch:s1}

\sessioncard{شنبه ۴ مهر}
{مفاهیم اولیهٔ احتمال و زمینه‌های شکل‌گیری؛ مقدمات: نظریهٔ مجموعه‌ها و ترکیبیات}
{\lr{Pishro-Nik} بخش‌های \lr{1.1} تا \lr{1.2} و \lr{2.1} \,؛\, \lr{Ross} فصل ۱ \,؛\, \lr{Papoulis} فصل ۱}
{\begin{itemize}
  \item توضیح دهید «احتمال» چه چیزی را مدل می‌کند و تعبیرهای کلاسیک، فراوانی نسبی و ذهنی چه تفاوتی دارند؛
  \item آزمایش تصادفی، فضای نمونه و پیشامد را تشخیص دهید و با زبان مجموعه‌ها بنویسید؛
  \item با اجتماع، اشتراک، متمم، افراز و قوانین دمورگان راحت کار کنید؛
  \item مجموعه‌های شمارا و ناشمارا را از هم تشخیص دهید و دلیلِ اهمیت این تمایز در احتمال را بدانید؛
  \item با اصل ضرب، جایگشت، ترکیب، ضرایب چندجمله‌ای و روش «ستاره و خط» مسائل شمارش را حل کنید.
\end{itemize}}

%% ------------------------------------------------------------
\section{احتمال چیست و چرا به آن نیاز داریم؟}
%% ------------------------------------------------------------

جهان پیرامون ما پر از \emph{عدم قطعیت} است. نمی‌دانیم فردا باران می‌بارد یا نه، نمی‌دانیم بستهٔ ارسالی در شبکه بدون خطا به مقصد می‌رسد یا نه، و نمی‌دانیم عمر یک قطعهٔ الکترونیکی چند ساعت خواهد بود. با این حال، در میانِ این بی‌نظمی \emph{نظم آماری} وجود دارد: اگر یک سکهٔ سالم را هزاران بار پرتاب کنیم، تقریباً نیمی از نتایج «شیر» خواهد بود. نظریهٔ احتمال، \key{ریاضیاتِ عدم قطعیت} است؛ زبانی دقیق برای توصیف، اندازه‌گیری و استدلال دربارهٔ پدیده‌هایی که نتیجه‌شان از پیش معلوم نیست.

چند نمونه از کاربردهای مهندسی:
\begin{itemize}
  \item \textbf{مخابرات:} کانال مخابراتی نویزی است؛ بیتِ ارسالیِ $0$ ممکن است $1$ دریافت شود. طراحی کدهای تصحیح خطا بدون احتمال ممکن نیست.
  \item \textbf{یادگیری ماشین و هوش مصنوعی:} تقریباً همهٔ مدل‌های یادگیری (از طبقه‌بند بیز ساده تا مدل‌های زبانی بزرگ) در اصل مدل‌های احتمالاتی‌اند.
  \item \textbf{شبکه و سیستم‌ها:} زمان رسیدن درخواست‌ها به سرور، طول صف‌ها و تأخیرها تصادفی‌اند.
  \item \textbf{الگوریتم‌ها:} الگوریتم‌های تصادفی (مثل \lr{Quicksort} تصادفی یا آزمون اول بودن میلر--رابین) با استفاده از تصادف سریع‌تر یا ساده‌تر می‌شوند.
  \item \textbf{قابلیت اطمینان:} احتمال خرابی یک سیستم از روی احتمال خرابی اجزای آن.
\end{itemize}

\subsection{دو روی یک سکه: نظریهٔ احتمال و آمار}

در این درس با دو نوع مسئلهٔ «معکوسِ یکدیگر» سروکار داریم:

\begin{center}
\begin{tikzpicture}[
  box/.style={draw=primary, fill=primarylight, rounded corners=4pt, minimum width=3.4cm,
              minimum height=1.1cm, align=center, font=\small},
  arr/.style={-{Stealth[length=3mm]}, line width=1pt}]
  \node[box] (m) at (0,0) {\rl{مدل احتمالاتی}\\ \lr{$(\Omega,\mathcal F,\mathbb P)$}};
  \node[box] (d) at (8,0) {\rl{داده‌ها و مشاهدات}};
  \draw[arr, defcolor] (m.north east) to[bend left=18]
     node[above, font=\small]{\rl{نظریهٔ احتمال: پیش‌بینیِ رفتار داده‌ها از روی مدل}} (d.north west);
  \draw[arr, excolor] (d.south west) to[bend left=18]
     node[below, font=\small]{\rl{آمار (استنتاج): حدس زدنِ مدل از روی داده‌ها}} (m.south east);
\end{tikzpicture}
\end{center}

\begin{itemize}
  \item \textbf{مدل‌سازی و نظریهٔ احتمال:} مدل معلوم است (مثلاً «سکه سالم است») و می‌پرسیم: احتمال اینکه در ۱۰ پرتاب دست‌کم ۸ شیر بیاید چقدر است؟ \quad (جلسه‌های ۱ تا ۲۱ این درس)
  \item \textbf{استنتاج آماری:} داده معلوم است (مثلاً «در ۱۰ پرتاب، ۸ شیر آمد») و می‌پرسیم: آیا سکه سالم است؟ احتمال شیر آمدن را چقدر تخمین بزنیم؟ \quad (جلسه‌های ۲۲ تا ۳۰)
\end{itemize}

\begin{history}{از میز قمار تا اصول موضوعهٔ کولموگروف}
\begin{itemize}[leftmargin=1.2em]
  \item \textbf{قرن ۱۶:} جرولامو کاردانو (\lr{Cardano}) در کتاب «دربارهٔ بازی‌های شانسی» نخستین محاسبات منظمِ احتمال را برای تاس انجام داد (انتشار پس از مرگ، ۱۶۶۳).
  \item \textbf{۱۶۵۴:} نامه‌نگاری مشهور \lr{Blaise Pascal} و \lr{Pierre de Fermat} دربارهٔ «مسئلهٔ تقسیم جایزه» و پرسش‌های شوالیه دومِره (\lr{de Méré}) دربارهٔ تاس؛ این سال را معمولاً \emph{تولد نظریهٔ احتمال} می‌دانند.
  \item \textbf{۱۶۵۷:} کریستیان هویگنس نخستین کتاب درسی احتمال را نوشت و مفهوم \emph{امید ریاضی} را معرفی کرد.
  \item \textbf{۱۷۱۳:} کتاب «هنر حدس زدن» یاکوب برنولی با نخستین نسخهٔ \emph{قانون اعداد بزرگ} منتشر شد.
  \item \textbf{۱۷۳۳ و ۱۷۶۳:} دموآور تقریب نرمال برای توزیع دوجمله‌ای را یافت؛ مقالهٔ توماس بیز (پس از مرگش) دربارهٔ «احتمال معکوس» منتشر شد.
  \item \textbf{۱۸۱۲:} لاپلاس در «نظریهٔ تحلیلی احتمالات» تعریف کلاسیک احتمال را صورت‌بندی کرد.
  \item \textbf{۱۹۳۳:} آندری کولموگروف در کتاب «مبانی نظریهٔ احتمال» احتمال را بر پایهٔ \emph{نظریهٔ مجموعه‌ها و نظریهٔ اندازه} بنا کرد و با سه اصل موضوع، آن را به شاخه‌ای دقیق از ریاضیات تبدیل کرد. (جلسهٔ ۲)
\end{itemize}
\end{history}

%% ------------------------------------------------------------
\section{تعبیرهای مختلف از احتمال}
%% ------------------------------------------------------------

وقتی می‌گوییم «احتمال شیر آمدن $\tfrac12$ است» دقیقاً چه منظوری داریم؟ در طول تاریخ چند پاسخ به این پرسش داده شده است.

\subsection{تعبیر کلاسیک (لاپلاس)}

\begin{definition}{احتمال کلاسیک}
اگر یک آزمایش $N$ نتیجهٔ ممکنِ \emph{هم‌شانس} داشته باشد و $N_A$ تا از آن‌ها «مساعدِ» رخداد $A$ باشند، آنگاه
\[ \P(A) = \frac{N_A}{N} = \frac{\text{تعداد حالت‌های مساعد}}{\text{تعداد کل حالت‌ها}}. \]
\end{definition}

این تعریف برای تاس و کارت و قرعه‌کشی بسیار کارآمد است، اما دو ایراد جدی دارد:
\begin{enumerate}
  \item \textbf{دوری است:} «هم‌شانس» یعنی «هم‌احتمال»؛ یعنی برای تعریف احتمال از خودِ احتمال استفاده کرده‌ایم!
  \item \textbf{محدود است:} برای تاسِ ناسالم، یا وقتی تعداد حالت‌ها نامتناهی است (مثلاً طول عمر یک لامپ)، کاربرد ندارد.
\end{enumerate}

\subsection{تعبیر فراوانی نسبی}

آزمایش را $N$ بار \emph{تحت شرایط یکسان} تکرار می‌کنیم و $N_A$ تعداد دفعاتی است که $A$ رخ داده است. نسبت $N_A/N$ را \term{فراوانی نسبی}{relative frequency} می‌نامیم و تعریف می‌کنیم:
\begin{keyformula}
\[ \P(A) = \lim_{N\to\infty} \frac{N_A}{N}, \qquad 0 \le \frac{N_A}{N} \le 1. \]
\end{keyformula}

\begin{example}{آزمایش‌های تاریخی پرتاب سکه}
چند آزمایش واقعی و مشهور:
\begin{center}
\begin{tabular}{lccc}
\toprule
آزمایشگر & تعداد پرتاب $N$ & تعداد شیر $N_A$ & فراوانی نسبی \\
\midrule
\lr{Buffon} (قرن ۱۸) & $4040$ & $2048$ & $0.5069$ \\
\lr{Kerrich} (۱۹۴۶) & $10000$ & $5067$ & $0.5067$ \\
\lr{Karl Pearson} (حدود ۱۹۰۰) & $24000$ & $12012$ & $0.5005$ \\
\bottomrule
\end{tabular}
\end{center}
با افزایش $N$ فراوانی نسبی به $0.5$ نزدیک‌تر می‌شود. این «پایداری فراوانی نسبی» پدیده‌ای تجربی است که بعدها (جلسهٔ ۲۱ به بعد) با \emph{قانون اعداد بزرگ} به صورت یک قضیهٔ ریاضی اثبات خواهیم کرد.
\end{example}

ایرادهای این تعبیر:
\begin{itemize}
  \item \textbf{وجود حد تضمین‌شده نیست:} از کجا بدانیم دنبالهٔ $N_A/N$ همگراست؟ این یک فرض تجربی است، نه یک حکم ریاضی.
  \item \textbf{همهٔ آزمایش‌ها تکرارپذیر نیستند:} «احتمال اینکه این پروژه تا پایان سال تمام شود» را نمی‌توان با هزار بار تکرار آزمایش سنجید.
  \item عملاً هرگز $N=\infty$ نداریم؛ پس همیشه با \emph{تخمین} سروکار داریم (موضوع بخش آمار درس).
\end{itemize}

\subsection{تعبیر ذهنی (بیزی)}
احتمال، \emph{درجهٔ باورِ} یک فرد به درستی یک گزاره است؛ مثلاً «به احتمال ۷۰٪ فردا باران می‌بارد». این تعبیر برای رخدادهای تکرارناپذیر معنا دارد و پایهٔ \emph{آمار بیزی} است (جلسه‌های ۲۶ و ۲۷). ایرادش این است که دو نفر ممکن است احتمال‌های متفاوتی به یک رخداد نسبت دهند.

\subsection{رویکرد اصل موضوعی (کولموگروف)}

کولموگروف به جای پاسخ به پرسش فلسفیِ «احتمال واقعاً چیست؟»، پرسید: \emph{احتمال، هر تعبیری که داشته باشد، باید از چه قواعدی پیروی کند؟} او سه اصل موضوع ساده وضع کرد و همهٔ نظریه را از آن‌ها نتیجه گرفت. هر سه تعبیرِ بالا این اصول را برآورده می‌کنند؛ پس نتایج نظریه برای همهٔ آن‌ها معتبر است.

\begin{intuition}
این رویکرد شبیه هندسهٔ اقلیدسی است: نمی‌پرسیم «نقطه واقعاً چیست؟»؛ فقط قواعد کار با نقطه و خط را می‌پذیریم و قضیه‌ها را اثبات می‌کنیم. برای این کار دو ابزار ریاضی لازم داریم:
(۱) \textbf{نظریهٔ مجموعه‌ها} تا رخدادها را به صورت مجموعه بنویسیم و با آن‌ها محاسبه کنیم؛
(۲) \textbf{نظریهٔ اندازه}، چون احتمال در واقع «اندازه‌ای» است که به مجموعه‌ها نسبت داده می‌شود، درست مثل طول، مساحت یا جرم.
\end{intuition}

%% ------------------------------------------------------------
\section{آزمایش تصادفی، فضای نمونه و پیشامد}
%% ------------------------------------------------------------

\begin{definition}{آزمایش تصادفی، برآمد، فضای نمونه}
\begin{itemize}
  \item \term{آزمایش تصادفی}{random experiment} فرایندی است که نتیجه‌اش از پیش به طور قطعی معلوم نیست، ولی مجموعهٔ همهٔ نتیجه‌های ممکن آن معلوم است.
  \item هر نتیجهٔ ممکنِ آزمایش را یک \term{برآمد}{outcome} می‌نامیم و معمولاً با $\omega$ نشان می‌دهیم.
  \item مجموعهٔ همهٔ برآمدهای ممکن را \term{فضای نمونه}{sample space} می‌نامیم و با $\Omega$ (یا $S$) نشان می‌دهیم.
\end{itemize}
\end{definition}

\begin{definition}{پیشامد}
هر زیرمجموعهٔ $A\subseteq \Omega$ (از زیرمجموعه‌های «مجاز»؛ جلسهٔ ۲ را ببینید) یک \term{پیشامد}{event} نام دارد. اگر برآمدِ آزمایش $\omega$ باشد و $\omega\in A$، می‌گوییم پیشامد $A$ \emph{رخ داده است}. پیشامدی که فقط یک عضو دارد، $\{\omega\}$، \term{پیشامد ساده}{elementary event} نامیده می‌شود.
\end{definition}

\begin{example}{چند فضای نمونه}
\begin{enumerate}
  \item پرتاب یک سکه: $\Omega=\{H,T\}$ \;(شیر $=H$، خط $=T$).
  \item پرتاب یک تاس: $\Omega=\{1,2,3,4,5,6\}$. پیشامد «عدد زوج بیاید» برابر $A=\{2,4,6\}$ است.
  \item پرتاب دو سکه (به ترتیب): $\Omega=\{HH,HT,TH,TT\}$. پیشامد «دقیقاً یک شیر» برابر $\{HT,TH\}$ است.
  \item تعداد ایمیل‌هایی که در یک ساعت می‌رسد: $\Omega=\{0,1,2,\dots\}$ (نامتناهیِ شمارا).
  \item طول عمر یک لامپ بر حسب ساعت: $\Omega=[0,\infty)$ (نامتناهیِ ناشمارا). پیشامد «لامپ بیش از ۱۰۰۰ ساعت کار کند» برابر $(1000,\infty)$ است.
\end{enumerate}
\end{example}

\begin{note}
انتخاب فضای نمونه به \emph{هدف ما} بستگی دارد. در پرتاب دو تاس، اگر فقط مجموع دو عدد برایمان مهم باشد، می‌توانیم $\Omega=\{2,\dots,12\}$ بگیریم؛ اما این برآمدها هم‌شانس نیستند. انتخاب $\Omega=\{1,\dots,6\}^2$ با ۳۶ برآمدِ هم‌شانس محاسبه را بسیار ساده‌تر می‌کند. قاعدهٔ کلی: فضای نمونه را آن‌قدر «ریز» بگیرید که همهٔ پرسش‌هایتان با آن قابل پاسخ باشد.
\end{note}

از اینجا به بعد، پیشامدها مجموعه‌اند و ترکیب کردن پیشامدها («$A$ یا $B$»، «$A$ و $B$»، «$A$ رخ ندهد») همان عملیات روی مجموعه‌هاست. پس ابتدا نظریهٔ مجموعه‌ها را مرور می‌کنیم.

%% ------------------------------------------------------------
\section{مروری بر نظریهٔ مجموعه‌ها}
%% ------------------------------------------------------------

\subsection{تعاریف پایه}

\begin{definition}{مجموعه، زیرمجموعه، تساوی}
\begin{itemize}
  \item \term{مجموعه}{set} گردایه‌ای از اشیای متمایز است که \emph{اعضای} آن نامیده می‌شوند. $x\in A$ یعنی $x$ عضو $A$ است و $x\notin A$ یعنی نیست.
  \item $A$ \term{زیرمجموعهٔ}{subset} $B$ است ($A\subseteq B$) اگر هر عضو $A$ عضو $B$ هم باشد.
  \item $A=B$ اگر و تنها اگر $A\subseteq B$ و $B\subseteq A$.
  \item \term{مجموعهٔ تهی}{empty set} $\varnothing$ هیچ عضوی ندارد و زیرمجموعهٔ هر مجموعه‌ای است.
  \item \term{مجموعهٔ مرجع}{universal set} مجموعه‌ای است که همهٔ مجموعه‌های مورد بحث زیرمجموعهٔ آن‌اند. در احتمال، مجموعهٔ مرجع همان $\Omega$ است.
\end{itemize}
\end{definition}

مجموعه‌ها را به دو شکل مشخص می‌کنیم: با \emph{فهرست کردن اعضا}، مثل $A=\{1,3,5\}$، یا با \emph{یک ویژگی}، مثل $A=\setc{x\in\N}{x<6}$ که همان $\{1,2,3,4,5\}$ است. مجموعه‌های عددی مهم: $\N=\{1,2,3,\dots\}$، $\Z$ (اعداد صحیح)، $\Q$ (اعداد گویا) و $\R$ (اعداد حقیقی).

\begin{warning}
نماد $\in$ برای «عضویت» و نماد $\subseteq$ برای «زیرمجموعه بودن» است. مثلاً برای $\Omega=\{1,\dots,6\}$ داریم $3\in\Omega$ و $\{3\}\subseteq\Omega$، ولی نوشتن $\{3\}\in\Omega$ یا $3\subseteq\Omega$ غلط است. در احتمال، $\omega$ یک \emph{برآمد} است ولی $\{\omega\}$ یک \emph{پیشامد}.
\end{warning}

\subsection{عملیات روی مجموعه‌ها}

\begin{definition}{اجتماع، اشتراک، متمم، تفاضل}
برای مجموعه‌های $A,B\subseteq\Omega$:
\begin{align*}
  A\cup B &= \setc{x}{x\in A \ \text{یا}\ x\in B} && \text{(اجتماع: «$A$ یا $B$»)}\\
  A\cap B &= \setc{x}{x\in A \ \text{و}\ x\in B} && \text{(اشتراک: «$A$ و $B$»)}\\
  A^c &= \setc{x\in\Omega}{x\notin A} && \text{(متمم: «$A$ رخ ندهد»)}\\
  A- B &= A\cap B^c && \text{(تفاضل: «$A$ رخ دهد ولی $B$ نه»)}\\
  A\,\triangle\, B &= (A- B)\cup(B- A) && \text{(تفاضل متقارن: «دقیقاً یکی»)}
\end{align*}
اشتراک $A\cap B$ را گاهی به اختصار $AB$ می‌نویسیم.
\end{definition}

\begin{figblock}
\begin{tikzpicture}[scale=0.62]
  \def\univ{(-1.7,-1.35) rectangle (2.9,1.35)}
  \def\sA{(0,0) circle (1)}
  \def\sB{(1.2,0) circle (1)}
  \tikzset{lbl/.style={font=\small}, venn/.style={draw=primary, thick}}
  % 1 union
  \begin{scope}
    \fill[vennfill] \sA; \fill[vennfill] \sB;
    \draw[venn] \sA \sB; \draw[graytext] \univ;
    \node at (-0.4,0.25) {$A$}; \node at (1.6,0.25) {$B$};
    \node[lbl] at (0.6,-1.85) {$A\cup B$};
  \end{scope}
  % 2 intersection
  \begin{scope}[xshift=5.4cm]
    \begin{scope}\clip \sA; \fill[vennfill] \sB; \end{scope}
    \draw[venn] \sA \sB; \draw[graytext] \univ;
    \node at (-0.4,0.25) {$A$}; \node at (1.6,0.25) {$B$};
    \node[lbl] at (0.6,-1.85) {$A\cap B$};
  \end{scope}
  % 3 complement
  \begin{scope}[xshift=10.8cm]
    \fill[vennfill] \univ; \fill[white] \sA;
    \draw[venn] \sA; \draw[graytext] \univ;
    \node at (0,0) {$A$};
    \node[lbl] at (0.6,-1.85) {$A^c$};
  \end{scope}
  % 4 difference
  \begin{scope}[yshift=-4.2cm]
    \fill[vennfill] \sA; \fill[white] \sB;
    \draw[venn] \sA \sB; \draw[graytext] \univ;
    \node at (-0.4,0.25) {$A$}; \node at (1.6,0.25) {$B$};
    \node[lbl] at (0.6,-1.85) {$A-B$};
  \end{scope}
  % 5 symmetric difference
  \begin{scope}[xshift=5.4cm, yshift=-4.2cm]
    \fill[vennfill, even odd rule] \sA \sB;
    \draw[venn] \sA \sB; \draw[graytext] \univ;
    \node at (-0.4,0.25) {$A$}; \node at (1.6,0.25) {$B$};
    \node[lbl] at (0.6,-1.85) {$A\,\triangle\,B$};
  \end{scope}
  % 6 disjoint
  \begin{scope}[xshift=10.8cm, yshift=-4.2cm]
    \draw[venn, fill=vennfill!80] (-0.5,0) circle (0.75);
    \draw[venn, fill=setB!35] (1.7,0) circle (0.75);
    \draw[graytext] \univ;
    \node at (-0.5,0) {$A$}; \node at (1.7,0) {$B$};
    \node[lbl] at (0.6,-1.85) {$A\cap B=\varnothing$};
  \end{scope}
\end{tikzpicture}

\figcaption[fig:venn-ops]{نمودارهای وِن برای عملیات مجموعه‌ای (ناحیهٔ رنگی نتیجهٔ عمل است).}
\end{figblock}

اجتماع و اشتراک را برای خانواده‌های بزرگ‌تر هم تعریف می‌کنیم:
\[
\bigcup_{i=1}^{\infty} A_i = \setc{x}{\exists\, i:\ x\in A_i},\qquad
\bigcap_{i=1}^{\infty} A_i = \setc{x}{\forall\, i:\ x\in A_i}.
\]
یعنی اجتماع، نقاطی است که \emph{دست‌کم در یکی} از $A_i$ها باشند و اشتراک، نقاطی که \emph{در همهٔ} $A_i$ها باشند.

\begin{example}{اجتماع و اشتراک نامتناهی}
فرض کنید $A_n=\left[0,\,1-\tfrac1n\right]$ و $B_n=\left(-\tfrac1n,\,\tfrac1n\right)$. آنگاه
\[ \bigcup_{n=1}^\infty A_n = [0,1), \qquad \bigcap_{n=1}^\infty B_n=\{0\}. \]
\begin{solution}
برای اولی: اگر $0\le x<1$ آنگاه $n$ای هست که $1-\tfrac1n\ge x$، پس $x$ در آن $A_n$ است. از طرف دیگر $1$ در هیچ $A_n$ای نیست، چون همواره $1-\tfrac1n<1$. برای دومی: $0$ در همهٔ $B_n$ها هست؛ و اگر $x\ne0$، برای $n>1/|x|$ داریم $x\notin B_n$.
توجه کنید که با آنکه همهٔ $A_n$ها \emph{بسته}‌اند، اجتماعشان بسته نیست. در جلسهٔ ۲ از چنین دنباله‌هایی برای «پیوستگی احتمال» استفاده می‌کنیم.
\end{solution}
\end{example}

\begin{theorem}[label=thm:set-algebra]{جبر مجموعه‌ها}
برای هر سه مجموعهٔ $A,B,C\subseteq\Omega$:
\begin{enumerate}
  \item \textbf{جابه‌جایی:} $A\cup B=B\cup A$ \quad و \quad $A\cap B=B\cap A$.
  \item \textbf{شرکت‌پذیری:} $A\cup(B\cup C)=(A\cup B)\cup C$ \quad و \quad $A\cap(B\cap C)=(A\cap B)\cap C$.
  \item \textbf{توزیع‌پذیری:} $A\cap(B\cup C)=(A\cap B)\cup(A\cap C)$ \quad و \quad $A\cup(B\cap C)=(A\cup B)\cap(A\cup C)$.
  \item \textbf{متمم:} $A\cup A^c=\Omega$، \ $A\cap A^c=\varnothing$، \ $(A^c)^c=A$، \ $\Omega^c=\varnothing$.
\end{enumerate}
\end{theorem}

\begin{proof}
همه با روش «شمول دوطرفه» ثابت می‌شوند. به عنوان نمونه، اولین قانون توزیع‌پذیری:
(نماد $\wedge$ یعنی «و» و $\vee$ یعنی «یا»)
\begin{align*}
x\in A\cap(B\cup C) &\iff x\in A \ \wedge\ (x\in B \ \vee\ x\in C)\\
&\iff (x\in A \ \wedge\ x\in B)\ \vee\ (x\in A \ \wedge\ x\in C)
\iff x\in (A\cap B)\cup(A\cap C).
\end{align*}
گام میانی همان توزیع‌پذیریِ «و» روی «یا» در منطق است.
\end{proof}

\begin{theorem}[label=thm:demorgan]{قوانین دمورگان}
برای هر خانوادهٔ (متناهی یا نامتناهی) از مجموعه‌های $A_1,A_2,\dots$:
\[
\Bigl(\bigcup_{i} A_i\Bigr)^c = \bigcap_{i} A_i^c,
\qquad
\Bigl(\bigcap_{i} A_i\Bigr)^c = \bigcup_{i} A_i^c.
\]
به‌ویژه $(A\cup B)^c=A^c\cap B^c$ و $(A\cap B)^c=A^c\cup B^c$.
\end{theorem}

\begin{proof}
برای قانون اول:
\[
x\in\Bigl(\bigcup_i A_i\Bigr)^c \iff \neg\bigl(\exists\, i:\ x\in A_i\bigr) \iff \forall i:\ x\notin A_i \iff \forall i:\ x\in A_i^c \iff x\in\bigcap_i A_i^c.
\]
به زبان ساده: $x$ در اجتماع نیست اگر و تنها اگر در \emph{هیچ} $A_i$ای نباشد.
قانون دوم از اعمال قانون اول به خانوادهٔ $\{A_i^c\}$ و گرفتن متمم از دو طرف به دست می‌آید:
$\bigl(\bigcup_i A_i^c\bigr)^c=\bigcap_i A_i$، پس $\bigcup_i A_i^c=\bigl(\bigcap_i A_i\bigr)^c$.
\end{proof}

\begin{intuition}
دمورگان به زبان پیشامدها: «\emph{نه} ($A$ یا $B$)» یعنی «($A$ نه) \emph{و} ($B$ نه)». مثلاً «هیچ‌کدام از دو سرور خراب نشود» یعنی «سرور اول خراب نشود و سرور دوم هم خراب نشود». این قانون در محاسبهٔ احتمال «دست‌کم یکی» بسیار کاربرد دارد: $\P(\text{دست‌کم یکی})=1-\P(\text{هیچ‌کدام})$.
\end{intuition}

\subsection{مجموعه‌های مجزا و افراز}

\begin{definition}{مجموعه‌های مجزا و افراز}
\begin{itemize}
  \item $A$ و $B$ \term{مجزا}{disjoint} (ناسازگار) هستند اگر $A\cap B=\varnothing$؛ یعنی دو پیشامد نمی‌توانند هم‌زمان رخ دهند.
  \item خانوادهٔ $A_1,A_2,\dots$ \emph{دوبه‌دو مجزا} است اگر به ازای هر $i\ne j$ داشته باشیم $A_i\cap A_j=\varnothing$.
  \item خانوادهٔ $B_1,B_2,\dots$ یک \term{افراز}{partition} برای $\Omega$ است اگر دوبه‌دو مجزا باشند و $\bigcup_i B_i=\Omega$. به بیان دیگر، \emph{دقیقاً یکی} از $B_i$ها رخ می‌دهد.
\end{itemize}
\end{definition}

\begin{figblock}
\begin{tikzpicture}[scale=0.9]
  \draw[graytext, thick, rounded corners=6pt] (0,0) rectangle (7,3.4);
  \draw[primary, thick] (2.1,0) -- (2.5,3.4);
  \draw[primary, thick] (4.3,0) -- (4.1,3.4);
  \draw[primary, thick] (4.3,1.5) -- (7,1.9);
  \fill[setB, opacity=0.35] (3.5,1.7) ellipse (2.2 and 1.05);
  \draw[setB!80!black, thick] (3.5,1.7) ellipse (2.2 and 1.05);
  \node at (0.7,2.9) {$B_1$}; \node at (3.2,3.05) {$B_2$};
  \node at (6.5,3.0) {$B_3$}; \node at (6.5,0.4) {$B_4$};
  \node[setB!60!black] at (3.4,1.6) {$A$};
  \node at (6.75,3.65) {$\Omega$};
\end{tikzpicture}

\figcaption[fig:partition]{افراز $\{B_1,\dots,B_4\}$ از $\Omega$؛ پیشامد $A$ به چهار تکهٔ مجزای $A\cap B_i$ شکسته می‌شود.}
\end{figblock}

\begin{theorem}[label=thm:partition-decomp]{شکستن یک مجموعه با افراز}
اگر $B_1,B_2,\dots$ افرازی از $\Omega$ باشد، آنگاه برای هر $A\subseteq\Omega$:
\[
A=\bigcup_i (A\cap B_i),
\]
و مجموعه‌های $A\cap B_i$ دوبه‌دو مجزا هستند. به‌ویژه برای هر $B$: \ $A=(A\cap B)\cup(A\cap B^c)$.
\end{theorem}

\begin{proof}
با توزیع‌پذیری (در حالت نامتناهی هم به همان روش): $\bigcup_i(A\cap B_i)=A\cap\bigl(\bigcup_i B_i\bigr)=A\cap\Omega=A$.
برای مجزا بودن: اگر $i\ne j$، آنگاه $(A\cap B_i)\cap(A\cap B_j)\subseteq B_i\cap B_j=\varnothing$.
حالت خاص از این نتیجه می‌شود که $\{B,B^c\}$ خودش یک افراز است.
\end{proof}

این قضیهٔ به‌ظاهر ساده، پایهٔ \emph{قانون احتمال کل} در جلسهٔ ۳ است.

\subsection{ضرب دکارتی و مجموعهٔ توانی}

\begin{definition}{ضرب دکارتی}
\term{ضرب دکارتی}{Cartesian product} دو مجموعه، مجموعهٔ همهٔ \emph{زوج‌های مرتب} است:
\[ A\times B=\setc{(a,b)}{a\in A,\ b\in B}. \]
به همین ترتیب $A_1\times\cdots\times A_n$ مجموعهٔ $n$-تایی‌های مرتب است و $A^n=A\times\cdots\times A$.
\end{definition}

\begin{example}{پرتاب دو تاس}
فضای نمونهٔ پرتاب دو تاس (اولی سفید و دومی قرمز) برابر است با
$\Omega=\{1,\dots,6\}\times\{1,\dots,6\}=\{1,\dots,6\}^2$ که ۳۶ عضو دارد. توجه کنید $(1,2)\ne(2,1)$: «سفید ۱ و قرمز ۲» با «سفید ۲ و قرمز ۱» فرق دارد.
\end{example}

\begin{definition}{مجموعهٔ توانی}
\term{مجموعهٔ توانی}{power set} یک مجموعهٔ $A$، مجموعهٔ \emph{همهٔ زیرمجموعه‌های} آن است و با $\mathcal P(A)$ یا $2^A$ نشان داده می‌شود.
\end{definition}

\begin{example}{}
برای $\Omega=\{H,T\}$:\quad $2^\Omega=\bigl\{\varnothing,\{H\},\{T\},\{H,T\}\bigr\}$. یعنی همهٔ پیشامدهای ممکن در پرتاب یک سکه این چهار پیشامدند: «هیچ‌چیز» (پیشامد ناممکن)، «شیر»، «خط» و «شیر یا خط» (پیشامد حتمی).
\end{example}

\begin{theorem}[label=thm:powerset]{اندازهٔ مجموعهٔ توانی}
اگر $|A|=n$ آنگاه $|2^A|=2^n$.
\end{theorem}

\begin{proof}
اعضای $A$ را $a_1,\dots,a_n$ بنامید. هر زیرمجموعهٔ $S\subseteq A$ را می‌توان با یک رشتهٔ دودویی $b_1b_2\cdots b_n$ نمایش داد که در آن $b_i=1$ اگر $a_i\in S$ و $b_i=0$ در غیر این صورت. این تناظر یک‌به‌یک و پوشاست (هر رشته دقیقاً یک زیرمجموعه را مشخص می‌کند). تعداد رشته‌های دودویی به طول $n$ برابر $2^n$ است (برای هر مکان دو انتخاب؛ اصل ضرب، بخش \ref{sec:counting}).
\end{proof}

\begin{note}
نماد $2^A$ از همین قضیه می‌آید. در جلسهٔ ۲ خواهیم دید که وقتی $\Omega$ متناهی یا شماراست، می‌توانیم \emph{همهٔ} زیرمجموعه‌ها را پیشامد بگیریم ($\F=2^\Omega$)؛ ولی برای $\Omega=\R$ این کار مشکل‌ساز است.
\end{note}

\subsection{تابع نشانگر}

\begin{definition}{تابع نشانگر}
\term{تابع نشانگرِ}{indicator function} مجموعهٔ $A\subseteq\Omega$ تابع $\ind_A:\Omega\to\{0,1\}$ است با
\[ \ind_A(\omega)=\begin{cases}1 & \omega\in A,\\ 0 & \omega\notin A.\end{cases} \]
\end{definition}

تابع نشانگر عملیات مجموعه‌ای را به عملیات جبری تبدیل می‌کند (اثبات با بررسی حالت‌ها ساده است):
\[
\ind_{A\cap B}=\ind_A\,\ind_B,\qquad
\ind_{A^c}=1-\ind_A,\qquad
\ind_{A\cup B}=\ind_A+\ind_B-\ind_A\ind_B .
\]
این ابزار بعدها در اثبات اصل شمول و عدم شمول و در تعریف متغیرهای تصادفی برنولی (جلسهٔ ۵) به کار می‌آید.

\subsection{کاردینالیتی: شمارا و ناشمارا}

تعداد اعضای مجموعهٔ $A$ را \term{کاردینالیتی}{cardinality} آن می‌نامیم و با $|A|$ نشان می‌دهیم. برای مجموعه‌های نامتناهی، «تعداد» را با تناظر یک‌به‌یک مقایسه می‌کنیم.

\begin{definition}{مجموعهٔ شمارا و ناشمارا}
\begin{itemize}
  \item $A$ \term{متناهی}{finite} است اگر $|A|<\infty$.
  \item $A$ \term{نامتناهیِ شمارا}{countably infinite} است اگر تناظری یک‌به‌یک و پوشا بین $A$ و $\N$ وجود داشته باشد؛ یعنی بتوان اعضای $A$ را به صورت یک \emph{دنباله} $a_1,a_2,a_3,\dots$ فهرست کرد.
  \item $A$ \term{شمارا}{countable} است اگر متناهی یا نامتناهیِ شمارا باشد؛ در غیر این صورت \term{ناشمارا}{uncountable} است.
\end{itemize}
\end{definition}

\begin{example}{اعداد صحیح شمارا هستند}
$\Z$ را می‌توان به صورت دنبالهٔ $0,\,1,\,-1,\,2,\,-2,\,3,\,-3,\dots$ فهرست کرد؛ پس $\Z$ شماراست. جالب است که $\N\subsetneq\Z$ ولی «به یک اندازه» عضو دارند!
\end{example}

\begin{theorem}[label=thm:NxN]{شمارا بودن $\N\times\N$ و $\Q$}
مجموعه‌های $\N\times\N$ و $\Q$ شمارا هستند.
\end{theorem}

\begin{proof}
زوج‌های $(i,j)$ را در یک جدول نامتناهی بچینید (سطر $i$، ستون $j$) و آن‌ها را روی قطرهای $i+j=2,3,4,\dots$ یکی‌یکی بشمارید (شکل \ref{fig:diagonal}):
\[ (1,1),\ (1,2),\ (2,1),\ (1,3),\ (2,2),\ (3,1),\ \dots \]
هر قطر تعداد متناهی عضو دارد، پس هر زوج بالاخره در جایگاهی متناهی از این فهرست ظاهر می‌شود. بنابراین $\N\times\N$ شماراست.

برای $\Q$: هر عدد گویای مثبت به صورت $a/b$ با $(a,b)\in\N\times\N$ نوشته می‌شود؛ با پیمایش همان جدول و \emph{رد کردن} کسرهای تکراری (مثل $2/4=1/2$)، فهرستی از همهٔ گویاهای مثبت به دست می‌آید. سپس مانند $\Z$، صفر و منفی‌ها را یکی در میان اضافه می‌کنیم.
\end{proof}

\begin{figblock}
\begin{tikzpicture}[scale=0.85]
  \foreach \i in {1,...,4} {
    \foreach \j in {1,...,4} {
      \node[font=\small] (n\i\j) at (1.3*\j,-0.9*\i) {$\frac{\j}{\i}$};
    }
  }
  \foreach \j in {1,...,4} \node[font=\scriptsize, graytext] at (1.3*\j,-0.25) {$a=\j$};
  \foreach \i in {1,...,4} \node[font=\scriptsize, graytext] at (0.35,-0.9*\i) {$b=\i$};
  \draw[-{Stealth}, accent, thick, rounded corners=3pt]
     (n11) -- (n12) -- (n21) -- (n31) -- (n22) -- (n13) -- (n14) -- (n23) -- (n32) -- (n41);
  \node[font=\small, graytext] at (6.8,-2.4) {$\cdots$};
  \node[font=\small, graytext] at (2.6,-4.3) {$\vdots$};
\end{tikzpicture}

\figcaption[fig:diagonal]{شمارش قطری کسرهای $a/b$ (ایدهٔ کانتور).}
\end{figblock}

\begin{theorem}{اجتماع شمارا از مجموعه‌های شمارا}
اگر $A_1,A_2,\dots$ همه شمارا باشند، آنگاه $\bigcup_{i=1}^\infty A_i$ هم شماراست.
\end{theorem}

\begin{proof}
اعضای $A_i$ را به صورت $a_{i1},a_{i2},\dots$ فهرست کنید. نگاشت $(i,j)\mapsto a_{ij}$ از $\N\times\N$ (یا زیرمجموعه‌ای از آن) \emph{پوشا} روی اجتماع است. چون $\N\times\N$ شماراست، با پیمایش قطری و رد کردن تکراری‌ها، همهٔ اعضای اجتماع را در یک دنباله فهرست می‌کنیم.
\end{proof}

\begin{theorem}[label=thm:cantor]{کانتور: بازهٔ $[0,1)$ ناشماراست}
مجموعهٔ اعداد حقیقی بازهٔ $[0,1)$ ناشماراست؛ در نتیجه $\R$ هم ناشماراست.
\end{theorem}

\begin{proof}
برهان خلف. فرض کنید $[0,1)$ شمارا باشد و فهرست $x_1,x_2,x_3,\dots$ همهٔ اعضایش را شامل شود. هر $x_i$ را به صورت اعشاری بنویسید (بدون دنبالهٔ بی‌پایان ۹):
\[
x_1=0.\,\underline{d_{11}}\,d_{12}\,d_{13}\cdots,\quad
x_2=0.\,d_{21}\,\underline{d_{22}}\,d_{23}\cdots,\quad
x_3=0.\,d_{31}\,d_{32}\,\underline{d_{33}}\cdots,\ \dots
\]
عدد $y=0.\,e_1e_2e_3\cdots$ را این‌گونه بسازید: اگر $d_{ii}\neq 5$ قرار دهید $e_i=5$ و اگر $d_{ii}=5$ قرار دهید $e_i=4$. آنگاه $y\in[0,1)$، ولی $y$ با هر $x_i$ دست‌کم در رقم $i$-ام فرق دارد؛ پس $y$ در فهرست نیست. این با فرض تناقض دارد.
\end{proof}

\begin{note}
چرا این تمایز در احتمال مهم است؟ اصل سوم کولموگروف (جلسهٔ ۲) جمع‌پذیری را فقط برای \emph{اجتماع‌های شمارا} می‌خواهد. اگر $\Omega$ شمارا باشد، احتمال هر پیشامد از جمعِ احتمال برآمدهایش به دست می‌آید. اما اگر $\Omega=[0,1]$ باشد، احتمال هر نقطهٔ منفرد صفر است و جمعِ ناشمارا تعریف نمی‌شود؛ به همین دلیل مدل‌های \emph{گسسته} (جلسهٔ ۲) و \emph{پیوسته} (جلسهٔ ۷) را جداگانه بررسی می‌کنیم.
\end{note}



\begin{example}{ترجمهٔ جمله به مجموعه}
سه قطعهٔ یک سیستم را بررسی می‌کنیم و $A_i$ پیشامدِ «قطعهٔ $i$ سالم است» باشد. پیشامدهای زیر را بنویسید:
(الف) همهٔ قطعات سالم‌اند؛ (ب) دست‌کم یک قطعه خراب است؛ (ج) دقیقاً یک قطعه سالم است.
\begin{solution}
(الف) $A_1\cap A_2\cap A_3$. \quad
(ب) $A_1^c\cup A_2^c\cup A_3^c$ که طبق دمورگان برابر $(A_1\cap A_2\cap A_3)^c$ است؛ یعنی متممِ (الف).\quad
(ج) $(A_1\cap A_2^c\cap A_3^c)\cup(A_1^c\cap A_2\cap A_3^c)\cup(A_1^c\cap A_2^c\cap A_3)$؛ سه مجموعهٔ این اجتماع دوبه‌دو مجزا هستند.
\end{solution}
\end{example}

%% ------------------------------------------------------------
\section{ترکیبیات: هنر شمردن}
\label{sec:counting}
%% ------------------------------------------------------------

در بسیاری از مسائل، برآمدها هم‌شانس‌اند و طبق تعبیر کلاسیک $\P(A)=|A|/|\Omega|$. پس محاسبهٔ احتمال به \emph{شمردن} تبدیل می‌شود. اما شمردن دقیق، به‌ویژه وقتی اعداد بزرگ‌اند، ابزار می‌خواهد.

\subsection{اصل جمع و اصل ضرب}

\begin{theorem}[label=thm:mult-principle]{اصل جمع و اصل ضرب}
\begin{enumerate}
  \item \textbf{اصل جمع:} اگر $A_1,\dots,A_m$ مجموعه‌های متناهیِ دوبه‌دو مجزا باشند، $\bigl|\bigcup_{i=1}^m A_i\bigr|=\sum_{i=1}^m |A_i|$.
  \item \textbf{اصل ضرب:} اگر یک کار در $r$ مرحله انجام شود، مرحلهٔ اول به $n_1$ روش، و به ازای هر انتخاب در مراحل قبل، مرحلهٔ $i$-ام به $n_i$ روش قابل انجام باشد، آنگاه کل کار به
  \[ n_1\times n_2\times\cdots\times n_r \]
  روش قابل انجام است. به‌ویژه $|A_1\times\cdots\times A_r|=|A_1|\cdots|A_r|$.
\end{enumerate}
\end{theorem}

\begin{proof}
اصل جمع همان تعریفِ شمردن است. برای اصل ضرب با $r=2$: حالت‌ها را بر اساس انتخاب مرحلهٔ اول به $n_1$ گروهِ مجزا تقسیم کنید؛ هر گروه $n_2$ عضو دارد، پس طبق اصل جمع کل حالت‌ها $\underbrace{n_2+\cdots+n_2}_{n_1}=n_1n_2$ است. حالت کلی با استقرا روی $r$ ثابت می‌شود: $r-1$ مرحلهٔ اول را یک «مرحلهٔ بزرگ» با $n_1\cdots n_{r-1}$ روش در نظر بگیرید.
\end{proof}

\begin{figblock}
\begin{tikzpicture}[
  level 1/.style={sibling distance=3.2cm, level distance=1.25cm},
  level 2/.style={sibling distance=1.05cm, level distance=1.25cm},
  every node/.style={font=\small},
  dot/.style={circle, fill=primary, inner sep=1.4pt},
  edge from parent/.style={draw=primary!60, thick}]
  \node[dot, label=above:{\rl{شروع}}] {}
    child { node[fill=defbg, draw=defcolor, rounded corners=2pt] {\lr{$a_1$}}
      child { node {\lr{$b_1$}} } child { node {\lr{$b_2$}} } child { node {\lr{$b_3$}} } }
    child { node[fill=defbg, draw=defcolor, rounded corners=2pt] {\lr{$a_2$}}
      child { node {\lr{$b_1$}} } child { node {\lr{$b_2$}} } child { node {\lr{$b_3$}} } };
  \node[font=\small, graytext, align=center] at (5.8,-1.25) {\rl{مرحلهٔ ۱: $n_1=2$ روش}};
  \node[font=\small, graytext, align=center] at (5.8,-2.5) {\rl{مرحلهٔ ۲: $n_2=3$ روش}};
  \node[font=\small, accent!80!black] at (5.8,-3.3) {\rl{کل: $2\times3=6$ مسیر}};
\end{tikzpicture}

\figcaption[fig:tree]{نمودار درختی اصل ضرب.}
\end{figblock}

\begin{example}{رمز عبور}
یک رمز عبور ۸ کاراکتری از حروف کوچک انگلیسی (۲۶ حرف) و ارقام (۱۰ رقم) ساخته می‌شود. (الف) چند رمز ممکن وجود دارد؟ (ب) چند رمز دست‌کم یک رقم دارند؟
\begin{solution}
(الف) هر کاراکتر $36$ انتخاب دارد؛ طبق اصل ضرب $36^8\approx 2.82\times10^{12}$.\\
(ب) شمردن مستقیم «دست‌کم یک رقم» سخت است؛ از \emph{متمم} استفاده می‌کنیم: رمزهای بدون رقم $26^8$ تا هستند، پس جواب $36^8-26^8\approx 2.61\times10^{12}$ است. (این ترفند «متمم» را زیاد به کار خواهیم برد.)
\end{solution}
\end{example}

\begin{example}{تعداد توابع}
تعداد توابع از یک مجموعهٔ $k$عضوی $\{x_1,\dots,x_k\}$ به یک مجموعهٔ $n$عضوی برابر $n^k$ است: برای تصویر هر $x_i$ مستقلاً $n$ انتخاب داریم. (قضیهٔ \ref{thm:powerset} حالت خاص $n=2$ است: هر زیرمجموعه یک تابع از $A$ به $\{0,1\}$ است.)
\end{example}

\subsection{چهار حالت نمونه‌گیری}

فرض کنید از یک ظرف با $n$ مهرهٔ \emph{متمایز}، $k$ مهره برمی‌داریم. دو پرسش تعیین‌کننده‌اند:
\begin{itemize}
  \item \textbf{ترتیب مهم است؟} آیا $(1,2)$ و $(2,1)$ دو نتیجهٔ متفاوت‌اند؟
  \item \textbf{با جایگذاری است؟} آیا مهرهٔ برداشته‌شده به ظرف برمی‌گردد (و ممکن است دوباره انتخاب شود)؟
\end{itemize}

\begin{center}
\renewcommand{\arraystretch}{1.6}
\begin{tabular}{>{\semifont}l|c|c}
 & \textbf{با جایگذاری} & \textbf{بدون جایگذاری} \\
\hline
\textbf{با ترتیب} & $n^k$ & $\displaystyle P^n_k=\frac{n!}{(n-k)!}$ \\
\hline
\textbf{بدون ترتیب} & $\displaystyle \binom{n+k-1}{k}$ & $\displaystyle \binom{n}{k}=\frac{n!}{k!\,(n-k)!}$ \\
\end{tabular}
\end{center}

در ادامه هر خانهٔ این جدول را اثبات می‌کنیم. خانهٔ «با ترتیب، با جایگذاری» مستقیماً از اصل ضرب به دست می‌آید.

\subsection{جایگشت}

\begin{definition}{جایگشت}
هر چینشِ مرتبِ $k$ شیء متمایز از میان $n$ شیء (بدون تکرار) را یک \term{جایگشتِ $k$تایی}{$k$-permutation} می‌نامیم. وقتی $k=n$، به آن \emph{جایگشت} $n$ شیء می‌گوییم. تعریف می‌کنیم $n!=n(n-1)\cdots 2\cdot 1$ و $0!=1$.
\end{definition}

\begin{theorem}{تعداد جایگشت‌ها}
تعداد جایگشت‌های $k$تایی از $n$ شیء متمایز برابر است با
\[ P^n_k = n(n-1)(n-2)\cdots(n-k+1)=\frac{n!}{(n-k)!}. \]
به‌ویژه $n$ شیء را به $n!$ طریق می‌توان مرتب کرد.
\end{theorem}

\begin{proof}
اصل ضرب: جایگاه اول $n$ انتخاب دارد، جایگاه دوم (پس از کنار گذاشتن شیء اول) $n-1$ انتخاب، \dots، و جایگاه $k$-ام $n-k+1$ انتخاب.
\end{proof}

\begin{example}{رتبه‌بندی و نشستن دور میز}
(الف) از ۱۰ تیم، به چند طریق می‌توان مقام‌های اول تا سوم را تعیین کرد؟ \\
(ب) ۶ نفر به چند طریق می‌توانند دور یک میزِ \emph{گرد} بنشینند (چینش‌هایی که با چرخش به هم تبدیل می‌شوند یکسان‌اند)؟
\begin{solution}
(الف) $P^{10}_3=10\cdot9\cdot8=720$.\\
(ب) اگر صندلی‌ها شماره داشتند، $6!=720$ چینش داشتیم. هر چینشِ دور میز با ۶ چرخش متناظر است که همه یکسان شمرده می‌شوند؛ پس جواب $6!/6=5!=120$ است. (به طور کلی $(n-1)!$.)
\end{solution}
\end{example}

\begin{note}
$n!$ بسیار سریع رشد می‌کند ($10!=3{,}628{,}800$ و $52!\approx8\times10^{67}$). برای تخمین آن از \term{فرمول استرلینگ}{Stirling's formula} استفاده می‌شود:
\[ n!\approx\sqrt{2\pi n}\,\Bigl(\frac{n}{e}\Bigr)^n. \]
مثلاً برای $n=10$ این فرمول $3{,}598{,}696$ می‌دهد که خطایش کمتر از یک درصد است.
\end{note}

\subsection{ترکیب و ضریب دوجمله‌ای}

\begin{definition}{ترکیب}
هر زیرمجموعهٔ $k$عضوی از یک مجموعهٔ $n$عضوی (بدون در نظر گرفتن ترتیب) را یک \term{ترکیبِ $k$تایی}{$k$-combination} می‌نامیم. تعداد آن‌ها را با $\binom{n}{k}$ (بخوانید «$n$ انتخاب $k$») نشان می‌دهیم و \term{ضریب دوجمله‌ای}{binomial coefficient} می‌نامیم. قرارداد: اگر $k<0$ یا $k>n$ آنگاه $\binom nk=0$.
\end{definition}

\begin{theorem}[label=thm:binom-formula]{فرمول ضریب دوجمله‌ای}
برای $0\le k\le n$:
\[ \binom{n}{k}=\frac{n!}{k!\,(n-k)!}. \]
\end{theorem}

\begin{proof}
جایگشت‌های $k$تایی را به دو روش می‌شماریم. روش اول: مستقیماً $\frac{n!}{(n-k)!}$ تا هستند. روش دوم (اصل ضرب در دو مرحله): ابتدا \emph{کدام} $k$ شیء را انتخاب کنیم ($\binom nk$ روش)، سپس آن‌ها را مرتب کنیم ($k!$ روش). پس
\[ \binom nk\cdot k! = \frac{n!}{(n-k)!} \quad\Longrightarrow\quad \binom nk=\frac{n!}{k!\,(n-k)!}. \]
\end{proof}

\begin{intuition}
این اثبات نمونه‌ای از تکنیک مهمِ \key{شمردن دوگانه} است: یک مجموعه را به دو روش می‌شماریم و دو جواب را برابر قرار می‌دهیم. هر ترکیب $k$تایی دقیقاً $k!$ بار در شمارش جایگشت‌ها «تکرار» شده است؛ پس بر $k!$ تقسیم می‌کنیم.
\end{intuition}

\begin{example}{تشکیل کمیته}
از ۵ زن و ۷ مرد، کمیته‌ای ۵ نفره با ۲ زن و ۳ مرد انتخاب می‌کنیم. (الف) چند کمیتهٔ مختلف ممکن است؟ (ب) اگر دو نفر از مردها حاضر به همکاری با هم نباشند چطور؟
\begin{solution}
(الف) طبق اصل ضرب: $\binom52\binom73=10\times35=350$.\\
(ب) انتخاب‌های «بدِ» مردها آن‌هایی‌اند که هر دو نفرِ ناسازگار را دارند: نفر سوم به $\binom51=5$ طریق انتخاب می‌شود. پس انتخاب‌های مجاز مردها $35-5=30$ و جواب $10\times30=300$ است.
\end{solution}
\end{example}

\begin{example}{مسیرهای شبکه‌ای}
در یک شبکهٔ شطرنجی، از نقطهٔ $(0,0)$ به $(m,n)$ فقط با گام‌های «یک واحد به راست» یا «یک واحد به بالا» می‌رویم. چند مسیر وجود دارد؟
\begin{solution}
هر مسیر دقیقاً $m+n$ گام دارد که $m$ تا از آن‌ها «راست» است. مسیر با انتخاب \emph{جایگاهِ} گام‌های راست کاملاً مشخص می‌شود؛ پس تعداد مسیرها $\binom{m+n}{m}$ است. مثلاً از $(0,0)$ به $(4,3)$: $\binom73=35$ مسیر.
\end{solution}
\end{example}

\begin{warning}
مسئله: از یک دست کارت ۵۲تایی، به چند طریق می‌توان ۳ کارت انتخاب کرد که \emph{دست‌کم یکی} آس باشد؟
استدلال غلط: «ابتدا یک آس انتخاب کن ($4$ روش)، سپس ۲ کارت دیگر از ۵۱ کارت باقی‌مانده ($\binom{51}{2}$ روش)؛ جواب $4\times1275=5100$.» این شمارش \emph{تکراری} دارد: دستِ $\{A\spadesuit, A\heartsuit, 7\clubsuit\}$ دو بار شمرده می‌شود (یک بار با $A\spadesuit$ به عنوان «آسِ اول» و یک بار با $A\heartsuit$).\\
روش درست با متمم: $\binom{52}{3}-\binom{48}{3}=22100-17296=4804$.
\end{warning}

\begin{theorem}[label=thm:binom-props]{خواص ضریب دوجمله‌ای}
برای اعداد صحیح $n\ge1$ و $0\le k\le n$:
\begin{enumerate}
  \item \textbf{تقارن:} $\binom nk=\binom n{n-k}$.
  \item \textbf{قاعدهٔ پاسکال:} $\binom nk=\binom{n-1}{k-1}+\binom{n-1}{k}$.
  \item \textbf{اتحاد واندرموند:} $\displaystyle\binom{m+n}{k}=\sum_{j=0}^{k}\binom mj\binom n{k-j}$.
\end{enumerate}
\end{theorem}

\begin{proof}
هر سه را با استدلال ترکیبیاتی (شمردن دوگانه) ثابت می‌کنیم.
\begin{enumerate}
  \item انتخاب $k$ عضو برای «ماندن» معادل انتخاب $n-k$ عضو برای «کنار گذاشتن» است.
  \item یک عضو خاص، مثلاً $a$، را در نظر بگیرید. زیرمجموعه‌های $k$عضوی دو دسته‌اند: آن‌هایی که $a$ را دارند (باید $k-1$ عضو دیگر از $n-1$ عضو باقی‌مانده انتخاب شود: $\binom{n-1}{k-1}$) و آن‌هایی که $a$ را ندارند (هر $k$ عضو از $n-1$ عضو دیگر: $\binom{n-1}{k}$). طبق اصل جمع، حکم ثابت است.
  \item گروهی متشکل از $m$ زن و $n$ مرد داریم و می‌خواهیم $k$ نفر انتخاب کنیم (سمت چپ). بر اساس تعداد زنانِ انتخاب‌شده، $j$، دسته‌بندی کنید: $\binom mj$ روش برای زنان و $\binom n{k-j}$ روش برای مردان. جمع روی $j$ سمت راست را می‌دهد.
\end{enumerate}
\end{proof}

قاعدهٔ پاسکال یعنی هر عدد در \emph{مثلث پاسکال} مجموع دو عدد بالای آن است:

\begin{figblock}
\begin{tikzpicture}[x=0.72cm, y=0.62cm]
  \foreach \n in {0,...,6} {
    \foreach \k in {0,...,\n} {
      \pgfmathtruncatemacro{\val}{round(\n!/(\k!*(\n-\k)!))}
      \node[font=\small, circle, inner sep=1.3pt, minimum size=1.5em,
            fill=primary!\fpeval{8+4*\n}] at (\k-\n/2, -\n) {\lr{$\val$}};
    }
    \node[font=\scriptsize, graytext] at (-4.8,-\n) {\lr{$n=\n$}};
  }
  \draw[accent, thick, -{Stealth}] (-0.5+0.25,-3+0.3) -- (0-0.15,-4+0.35);
  \draw[accent, thick, -{Stealth}] (0.5-0.25,-3+0.3) -- (0+0.15,-4+0.35);
\end{tikzpicture}

\figcaption[fig:pascal]{مثلث پاسکال؛ مثلاً $\binom42=\binom31+\binom32=3+3=6$.}
\end{figblock}

\begin{theorem}[label=thm:binomial]{قضیهٔ دوجمله‌ای}
برای هر $x,y\in\R$ و هر عدد صحیح $n\ge0$:
\[ (x+y)^n=\sum_{k=0}^{n}\binom nk x^k y^{n-k}. \]
\end{theorem}

\begin{proof}
$(x+y)^n=(x+y)(x+y)\cdots(x+y)$ را باز کنید. هر جمله از انتخاب $x$ یا $y$ از \emph{هر} پرانتز به دست می‌آید. جملهٔ $x^ky^{n-k}$ دقیقاً وقتی ساخته می‌شود که از $k$ پرانتز $x$ و از بقیه $y$ انتخاب شود؛ تعداد روش‌های انتخابِ این $k$ پرانتز $\binom nk$ است. پس ضریب $x^ky^{n-k}$ برابر $\binom nk$ است.
\end{proof}

\begin{corollary}[label=cor:binom-sums]{}
\[ \sum_{k=0}^n\binom nk=2^n, \qquad \sum_{k=0}^n(-1)^k\binom nk=0\quad(n\ge1). \]
\end{corollary}

\begin{proof}
در قضیهٔ دوجمله‌ای یک بار $x=y=1$ و یک بار $x=-1,\ y=1$ قرار دهید.
\end{proof}

\begin{intuition}
اتحاد اول دوباره قضیهٔ \ref{thm:powerset} است: زیرمجموعه‌های یک مجموعهٔ $n$عضوی را بر اساس اندازه‌شان دسته‌بندی کرده‌ایم. اتحاد دوم می‌گوید تعداد زیرمجموعه‌های با اندازهٔ زوج با تعداد زیرمجموعه‌های با اندازهٔ فرد برابر است. در جلسهٔ ۲ همین اتحادها ثابت می‌کنند احتمال‌های «$k$ شیر در $n$ پرتاب» جمعشان ۱ است.
\end{intuition}

\subsection{ضریب چندجمله‌ای: جایگشت با تکرار}

\begin{theorem}[label=thm:multinomial]{ضریب چندجمله‌ای}
تعداد روش‌های تقسیم $n$ شیء متمایز به $r$ گروهِ \emph{برچسب‌دار} با اندازه‌های $n_1,\dots,n_r$ (که $n_1+\cdots+n_r=n$) برابر است با
\[ \binom{n}{n_1,n_2,\dots,n_r}=\frac{n!}{n_1!\,n_2!\cdots n_r!}. \]
به طور معادل، این تعداد جایگشت‌های متمایزِ $n$ شیء است که $n_1$ تا از نوع ۱، \dots، و $n_r$ تا از نوع $r$ (و اشیای هم‌نوع یکسان) باشند.
\end{theorem}

\begin{proof}
اصل ضرب: ابتدا $n_1$ شیء برای گروه اول ($\binom{n}{n_1}$ روش)، سپس $n_2$ شیء از $n-n_1$ شیء باقی‌مانده برای گروه دوم ($\binom{n-n_1}{n_2}$ روش)، و الی آخر:
\[
\binom{n}{n_1}\binom{n-n_1}{n_2}\cdots\binom{n_r}{n_r}
=\frac{n!}{n_1!\,(n-n_1)!}\cdot\frac{(n-n_1)!}{n_2!\,(n-n_1-n_2)!}\cdots
=\frac{n!}{n_1!\,n_2!\cdots n_r!}.
\]
(جمله‌ها به صورت تلسکوپی ساده می‌شوند.)
\end{proof}

\begin{example}{حروف یک کلمه}
با حروف کلمهٔ \lr{STATISTICS} چند رشتهٔ ۱۰ حرفیِ متمایز می‌توان ساخت؟
\begin{solution}
حروف: $S$ سه بار، $T$ سه بار، $I$ دو بار، $A$ یک بار، $C$ یک بار. جواب
\[ \frac{10!}{3!\,3!\,2!\,1!\,1!}=\frac{3{,}628{,}800}{72}=50{,}400. \]
\end{solution}
\end{example}

\begin{example}{تیم‌بندی}
۱۲ دانشجو را به ۳ گروه ۴ نفره تقسیم می‌کنیم. (الف) اگر گروه‌ها نام داشته باشند («الف»، «ب»، «ج») چند تقسیم ممکن است؟ (ب) اگر گروه‌ها بی‌نام باشند چطور؟
\begin{solution}
(الف) $\dfrac{12!}{4!\,4!\,4!}=34{,}650$. \quad
(ب) هر تقسیمِ بی‌نام با $3!$ نام‌گذاری متناظر است؛ پس $34{,}650/6=5{,}775$.
\end{solution}
\end{example}

\begin{note}
\textbf{قضیهٔ چندجمله‌ای} تعمیم قضیهٔ دوجمله‌ای است و با همان استدلال ثابت می‌شود:
\[ (x_1+\cdots+x_r)^n=\sum_{n_1+\cdots+n_r=n}\binom{n}{n_1,\dots,n_r}x_1^{n_1}\cdots x_r^{n_r}. \]
\end{note}

\subsection{نمونه‌گیری بدون ترتیب با جایگذاری: ستاره و خط}

\begin{theorem}[label=thm:stars-bars]{ستاره و خط}
\begin{enumerate}
  \item تعداد جواب‌های معادلهٔ $x_1+x_2+\cdots+x_n=k$ در اعداد صحیحِ \emph{نامنفی} برابر $\binom{n+k-1}{k}=\binom{n+k-1}{n-1}$ است.
  \item به طور معادل، تعداد روش‌های انتخاب $k$ شیء از $n$ نوع، \emph{با تکرار و بدون ترتیب}، برابر $\binom{n+k-1}{k}$ است.
  \item تعداد جواب‌ها در اعداد صحیحِ \emph{مثبت} ($x_i\ge1$) برابر $\binom{k-1}{n-1}$ است.
\end{enumerate}
\end{theorem}

\begin{proof}
(۱) هر جواب را با رشته‌ای از $k$ ستاره ($\star$) و $n-1$ خط ($|$) نمایش دهید: $x_1$ ستاره، سپس یک خط، سپس $x_2$ ستاره، و الی آخر. مثلاً برای $n=4,\ k=6$:
\[ \star\star\,|\,\star\,|\,\,|\,\star\star\star \quad\longleftrightarrow\quad (x_1,x_2,x_3,x_4)=(2,1,0,3). \]
این تناظر یک‌به‌یک و پوشاست. رشته طول $n+k-1$ دارد و با انتخاب جایگاه $k$ ستاره مشخص می‌شود: $\binom{n+k-1}{k}$.\\
(۲) $x_i$ را «تعداد دفعاتی که نوع $i$ انتخاب شده» تعبیر کنید.\\
(۳) با تغییر متغیر $y_i=x_i-1\ge0$ معادله به $y_1+\cdots+y_n=k-n$ تبدیل می‌شود که طبق (۱) تعداد جواب‌هایش $\binom{n+(k-n)-1}{k-n}=\binom{k-1}{k-n}=\binom{k-1}{n-1}$ است.
\end{proof}

\begin{example}{تقسیم آبنبات}
۱۰ آبنبات \emph{یکسان} را بین ۴ کودک تقسیم می‌کنیم. (الف) به چند طریق؟ (ب) اگر هر کودک دست‌کم یک آبنبات بگیرد؟
\begin{solution}
(الف) تعداد جواب‌های نامنفیِ $x_1+x_2+x_3+x_4=10$: \ $\binom{13}{3}=286$. \\
(ب) تعداد جواب‌های مثبت: $\binom{9}{3}=84$.
\end{solution}
\end{example}

\begin{warning}
در شمارش، همیشه دو پرسش را صریحاً جواب دهید: «\emph{اشیا متمایزند یا یکسان؟}» و «\emph{ظرف‌ها (گروه‌ها) متمایزند یا یکسان؟}». اگر ۱۰ آبنبات \emph{متمایز} بودند، جواب مثال بالا $4^{10}$ بود (هر آبنبات مستقلاً یکی از ۴ کودک را انتخاب می‌کند)، نه $286$.
همچنین به خاطر بسپارید که در \emph{محاسبهٔ احتمال}، برآمدهای «ستاره و خط» معمولاً هم‌شانس \emph{نیستند}؛ این نکته را در جلسهٔ ۲ بیشتر بررسی می‌کنیم.
\end{warning}

%% ------------------------------------------------------------
\section*{جمع‌بندی جلسهٔ اول}
\addcontentsline{toc}{section}{جمع‌بندی و تمرین‌ها}
%% ------------------------------------------------------------

\begin{summary}
\begin{itemize}
  \item \textbf{احتمال} ریاضیاتِ عدم قطعیت است؛ \emph{نظریهٔ احتمال} از مدل به داده می‌رود و \emph{آمار} از داده به مدل.
  \item تعبیرهای کلاسیک ($N_A/N$ برای حالت‌های هم‌شانس)، فراوانی نسبی ($\lim N_A/N$) و ذهنی هر کدام محدودیت دارند؛ رویکرد \textbf{اصل موضوعی کولموگروف} همه را پوشش می‌دهد.
  \item \textbf{فضای نمونه} $\Omega$ مجموعهٔ همهٔ برآمدهاست و \textbf{پیشامد} زیرمجموعه‌ای از آن. «یا» $\leftrightarrow\cup$، «و» $\leftrightarrow\cap$، «نه» $\leftrightarrow{}^c$.
  \item قوانین \textbf{دمورگان}: $(\bigcup A_i)^c=\bigcap A_i^c$ و $(\bigcap A_i)^c=\bigcup A_i^c$.
  \item \textbf{افراز}: پیشامدهای دوبه‌دو مجزا که اجتماعشان $\Omega$ است؛ هر $A$ برابر $\bigcup_i(A\cap B_i)$ است.
  \item $\N,\Z,\Q$ و $\N\times\N$ \textbf{شمارا} هستند ولی $[0,1)$ و $\R$ \textbf{ناشمارا}.
  \item شمارش: اصل ضرب؛ $n^k$ (با ترتیب، با جایگذاری)، $\frac{n!}{(n-k)!}$ (با ترتیب، بدون جایگذاری)، $\binom nk$ (بدون ترتیب، بدون جایگذاری)، $\binom{n+k-1}{k}$ (بدون ترتیب، با جایگذاری)، و ضریب چندجمله‌ای $\frac{n!}{n_1!\cdots n_r!}$.
  \item ترفندهای کلیدی: \key{متمم} (برای «دست‌کم یکی»)، \key{شمردن دوگانه}، و پرهیز از \key{شمارش تکراری}.
\end{itemize}
\end{summary}


\subsection*{تمرین‌های پیشنهادی}

\begin{exercise}
برای مجموعه‌های دلخواه $A,B,C$ ثابت کنید:
(الف) $A-(B\cup C)=(A-B)\cap(A-C)$؛ \
(ب) $A\cup B=A\cup(B-A)$ و دو مجموعهٔ سمت راست مجزا هستند؛ \
(ج) $A\,\triangle\,B=(A\cup B)-(A\cap B)$.
\end{exercise}

\begin{exercise}
یک سکه را آن‌قدر پرتاب می‌کنیم تا اولین شیر بیاید. فضای نمونه را بنویسید. آیا شماراست؟ پیشامد «تعداد پرتاب‌ها زوج باشد» را به صورت مجموعه بنویسید.
\end{exercise}

\begin{exercise}
ثابت کنید مجموعهٔ همهٔ رشته‌های دودوییِ \emph{متناهی} شماراست، ولی مجموعهٔ همهٔ دنباله‌های دودوییِ \emph{نامتناهی} ناشماراست. (راهنمایی: برای دومی از استدلال قطری کانتور استفاده کنید.)
\end{exercise}

\begin{exercise}
یک پلاک از ۲ رقم، سپس یک حرف از ۳۲ حرف الفبای فارسی و سپس ۳ رقم تشکیل شده است. (الف) چند پلاک ممکن است؟ (ب) در چند پلاک هیچ رقمی تکرار نشده است؟
\answer{(الف) $10^5\times32=3{,}200{,}000$ \quad (ب) $32\times10\cdot9\cdot8\cdot7\cdot6=967{,}680$}
\end{exercise}

\begin{exercise}
با استدلال ترکیبیاتی ثابت کنید: (الف) $\displaystyle k\binom nk=n\binom{n-1}{k-1}$ \ (راهنمایی: انتخاب یک کمیته و رئیس آن)؛ \ (ب) $\displaystyle\sum_{k=0}^n\binom nk^2=\binom{2n}{n}$.
\end{exercise}

\begin{exercise}
معادلهٔ $x_1+x_2+x_3=12$ چند جواب صحیح دارد اگر (الف) $x_i\ge0$؛ (ب) $x_i\ge1$؛ (ج) $x_1\ge2$ و $x_2\ge0$ و $x_3\ge3$؟
\answer{(الف) $91$ \quad (ب) $55$ \quad (ج) $36$}
\end{exercise}

\begin{exercise}[\lr{Ross}]
۸ رخ روی صفحهٔ شطرنج ۸×۸ به چند طریق قرار می‌گیرند به طوری که هیچ دو رخی یکدیگر را تهدید نکنند (هیچ دو رخی در یک سطر یا یک ستون نباشند)؟
\answer{$8!=40{,}320$}
\end{exercise}

\begin{exercise}\hard
با استفاده از قضیهٔ دوجمله‌ای ثابت کنید $\displaystyle\sum_{k=0}^n k\binom nk=n\,2^{n-1}$. سپس همین اتحاد را با یک استدلال ترکیبیاتی (شمردن دوگانه) هم ثابت کنید.
\end{exercise}
